\documentclass[10pt,a4paper]{amsart}
\usepackage[margin=19mm]{geometry}
\usepackage{amsmath,amssymb,mathtools}
\usepackage[scr=rsfs,cal=euler]{mathalfa}
\usepackage{xfrac}
\usepackage[unicode,hidelinks,bookmarks=false]{hyperref}
\newcommand{\ie}{i.e.\ }
\newcommand{\cf}{cf.\ }
\newcommand{\resp}{respectively\ }
\newcommand{\Subsection}{\subsection}
\providecommand{\nobreakdash}{\nobreak\hbox{-}}
\setlength{\emergencystretch}{2em}
\multlinegap0pt
\usepackage{bm}
\usepackage{bbm}
\usepackage{dsfont}
\usepackage{xspace}
\usepackage{cancel}
\usepackage{mathtools}
\usepackage{amsthm}
\makeatletter
\@ifundefined{proposition}{\newtheorem{proposition}{Proposition}}{}
\@ifundefined{theorem}{\newtheorem{theorem}{Theorem}}{}
\makeatother
\newcommand{\tri}[3]{\ensuremath{#1^{#2} - #1^{#3} + 1}}
\title[Dyadically resolving trinomials for fast modular arithmetic]{Dyadically resolving trinomials for fast modular arithmetic}
\author{Robert Dougherty-Bliss, Mits Kobayashi, Natalya Ter-Saakov, Eugene Zima}
\date{Source: arXiv:2508.11043v1 (CC BY 4.0)}
\begin{document}
\maketitle

\noindent\emph{Source and licence.} Dyadically resolving trinomials for fast modular arithmetic. Version arXiv:2508.11043v1 (CC BY 4.0), distributed under the Creative Commons Attribution 4.0 International licence, as recorded in the arXiv licence metadata for this exact version and shown on the abstract page https://arxiv.org/abs/2508.11043v1. Full rights notice: licenses/arxiv-2508.11043-CC-BY-4.0.txt.\par\smallskip

\noindent\textit{Editorial scope.} Counted unit: the Theorem whose source label is \texttt{None} in the article (source0.tex lines 740--742, source label \texttt{None}), delivered together with the article's own complete original proof of it (source0.tex lines 744--781, one \texttt{proof} environment of 38 lines). No mathematics is written, repaired or re-derived: statement and proof are the article's own text line for line. Required context retained uncounted: same-primes (lines 664--668). Every definition, display and label of those lines is retained unchanged. Omitted from this package and not redistributed: 1--663, 669--739, 743--743, 782--811. Nothing inside a retained excerpt is omitted or altered: every retained body is delivered exactly as the article's lines are written, so no omission receipt of any kind accompanies this package. Citations: the retained excerpts carry no citation command at all, so no bibliography entry of the article is transcribed and no reference is redistributed. Reference adaptation: none was needed and none was made; every reference inside the delivered unit resolves inside it, so no unresolved reference or citation is produced. Printed slips kept literal: no printed word, symbol, hypothesis or number of the counted statement or of its proof was altered. Layout and class: the article's own class is \texttt{article} with packages including biblatex, amsmath, amsthm, amsfonts, geometry, tcolorbox, hyperref; the wrapper is the lane's pinned \texttt{amsart} editorial wrapper at 10pt on A4, which restates the theorem-like environments the retained text needs and adds the article's own macro definitions that the retained text needs (\texttt{\textbackslash tri}), with extra wrapper packages (bm, bbm, dsfont, xspace, cancel, mathtools). The delivered numbering is the unit's own. Assets: no retained span contains a figure environment, a graphics-inclusion command, a PDF-page inclusion, a table or a TikZ picture; no image file is copied into this package, the package declares no asset dependency and no image file is redistributed with it. Licence: version arXiv:2508.11043v1 of this article is distributed under the Creative Commons Attribution 4.0 International licence, as recorded in the arXiv licence metadata for this exact version and shown on the abstract page https://arxiv.org/abs/2508.11043v1. A full rights notice is delivered at licenses/arxiv-2508.11043-CC-BY-4.0.txt. Characters: every retained excerpt is pure ASCII, so the delivered engine, XeLaTeX with its default Latin Modern text font, reports no missing character.
\begin{proposition}
    \label{same-primes}
    If $f(x)$ and $g(x)$ are monic, then their reduced resultant and resultant
    have the same prime divisors.
\end{proposition}

\begin{theorem}
    The moduli $2^n - 2^k + 1$ and $2^n - 2^j + 1$ are scalable if and only if $\tri{x}{n}{k}$ and $\tri{x}{n}{j}$ dyadically resolve.
\end{theorem}

\begin{proof}
    Define the polynomials
    \begin{align*}
        f(x) &= x^n - x^k + 1 \\
        g(x) &= x^n - x^j + 1.
    \end{align*}
    If the moduli are scalable, then there is some inverse
    polynomial $A(x)$ with dyadic coefficients such that $A(2^c) f(2^c) \equiv
    1 \pmod{g(2^c)}$ for all sufficiently large $c$. The previous lemma implies
    that $A(x) f(x) - 1$ is a multiple of $g(x)$, say $A(x) f(x) - B(x) g(x) =
    1$. Further, because $A$, $f$, and $g$ have dyadic coefficients, so does
    $B$. By clearing denominators, we obtain a power of $2$ as a
    $\mathbb{Z}[x]$-linear combination of $f$ and $g$, which implies that their
    reduced resultant is a power of 2. By Proposition~\ref{same-primes}, $f$
    and $g$ dyadically resolve.

    On the other hand, if $f$ and $g$ dyadically resolve, then their reduced
    B\'ezout cofactors $a$ and $b$ in
    \begin{equation*}
        a(x) f(x) + b(x) g(x) = 1
    \end{equation*}
    can be constructed using only dyadic coefficients. These cofactors are
    essentially our inverse polynomials, but they need some minor adjustments to
    match our definition.

    Without loss of generality, we will describe how to make $a$ match the
    definition. If $a(0)$ is not an integer, then the equation $a(0) + b(0) = 1$
    implies $a(0) = c/2^m$ and $b(0) = 1 - c/2^m$ for some integers $c$ and $m >
    1$. Replacing $a$ with $a - gc / 2^m$ and $b$ with $b + fc / 2^m$ makes the
    constant terms integers (at the expense of having degree at most $n$ rather
    than $n - 1$). Now it is true that $a(2^c) f(2^c) \equiv 1 \pmod{g(2^c)}$
    for sufficiently large $c$. If the leading coefficient of $a$ is positive,
    then we are done. If the leading coefficient is negative, then we can
    replace $a$ with $a + c' g$ and $g$ with $g - c' f$ for some sufficiently
    large integer $c'$. Now $a$ satisfies the definition. (Note that the $b$
    which results from this process has a negative leading coefficient. It must
    be repeated in full on $b$ to get a suitable inverse polynomial.)
\end{proof}

\end{document}
